% Part: first-order-logic
% Chapter: natural-deduction
% Section: quantifier-rules

\documentclass[../../../include/open-logic-section]{subfiles}

\begin{document}

\olfileid{fol}{ntd}{qrl}

\olsection{Kwantorregels}

\subsection{Regels voor $\lforall$}

\begin{defish}
\AxiomC{$!A(a)$}
\RightLabel{\Intro{\lforall}}
\UnaryInfC{$\lforall[x][\Atom{!A}{x}]$}
\DisplayProof
\hfill
\AxiomC{$\lforall[x][\Atom{!A}{x}]$}
\RightLabel{\Elim{\lforall}}
\UnaryInfC{$!A(t)$}
\DisplayProof
\end{defish}

In de regels voor~$\lforall$ is $t$ een gesloten term (een term die
geen variabelen bevat), en is $a$~!!a{constant} die niet voorkomt in
de conclusie~$\lforall[x][!A(x)]$ of in enige aanname die
!!{undischarged} is in de !!{derivation} die eindigt met de
premisse~$!A(a)$. We noemen $a$ de \emph{eigenvariabele} van de
afleidingsstap \Intro{\lforall}.\footnote{We gebruiken de term
``eigenvariabele'', hoewel $a$ in de bovenstaande regel een constante
is. Dit heeft historische redenen.}

\subsection{Regels voor $\lexists$}

\begin{defish}
\AxiomC{$\Atom{!A}{t}$}
\RightLabel{\Intro{\lexists}}
\UnaryInfC{$\lexists[x][\Atom{!A}{x}]$}
\DisplayProof
\hfill
\AxiomC{$\lexists[x][\Atom{!A}{x}]$}
\AxiomC{[$\Atom{!A}{a}$]$^n$}
\DeduceC{$!C$}
\DischargeRule{\Elim{\lexists}}{n}
\BinaryInfC{$!C$}
\DisplayProof
\end{defish}

Ook hier is $t$ een gesloten term, en is $a$ !!a{constant} die niet
voorkomt in de premisse $\lexists[x][!A(x)]$, in de conclusie~$!C$ of
in enige aanname die !!{undischarged} is in de !!{derivation}s die
eindigen met de twee premissen (afgezien van de aannamen $!A(a)$). We
noemen $a$ de \emph{eigenvariabele} van de afleidingsstap
\Elim{\lexists}.

De voorwaarde dat een eigenvariabele noch in de premissen voorkomt,
noch in enige aanname die !!{undischarged} is in de !!{derivation}s
die leiden tot de premissen van de afleidingsstap \Intro{\lforall} of
\Elim{\lexists}, heet de \emph{eigenvariabelevoorwaarde}.

\begin{explain}
Herinner u de afspraak dat we, wanneer $!A$ !!a{formula} is waarin de
!!{variable}~$x$ vrij voorkomt, dit aangeven door~$!A(x)$ te schrijven.
In dezelfde context is $!A(t)$ dan een afkorting voor
$\Subst{!A}{t}{x}$. We zouden de regel $\Intro\lexists$ dus ook als
volgt kunnen schrijven:
\begin{prooftree}
\AxiomC{$\Subst{!A}{t}{x}$}
\RightLabel{\Intro{\lexists}}
\UnaryInfC{$\lexists[x][!A]$}
\end{prooftree}
Merk op dat $t$ al in~$!A$ kan voorkomen; zo kan $!A$ bijvoorbeeld
$\Atom{\Obj P}{t,x}$ zijn. Uit~$\Atom{\Obj P}{t,t}$
$\lexists[x][\Atom{\Obj P}{t,x}]$ afleiden is dus een correcte
toepassing van~$\Intro\lexists$---men mag één of meer, en niet
noodzakelijk alle, voorkomens van~$t$ in de premisse door de gebonden
!!{variable}~$x$ ``vervangen''. De eigenvariabelevoorwaarden in
$\Intro\lforall$ en~$\Elim\lexists$ vereisen echter dat de
!!{constant}~$a$ niet in~$!A$ voorkomt. Men kan dus niet met een
correcte toepassing van~$\Intro\lforall$ uit $\Atom{\Obj P}{a,a}$ de
formule $\lforall[x][\Atom{\Obj P}{a,x}]$ afleiden.
\end{explain}

\begin{explain}
Voor \Intro{\lexists} en \Elim{\lforall} gelden geen beperkingen: de
term~$t$ mag willekeurig zijn, zodat we daar geen voorwaarden hoeven te
controleren. Bij de regels \Elim{\lexists} en \Intro{\lforall} vereist
de eigenvariabelevoorwaarde daarentegen dat de !!{constant}~$a$ nergens
in de conclusie of in een !!{undischarged} aanname voorkomt. Deze
voorwaarde is nodig om te waarborgen dat het systeem correct is, dat
wil zeggen dat het alleen !!{sentence}s !!{derive} uit
!!{undischarged} aannamen waaruit zij volgen. Zonder deze voorwaarde
zou het volgende zijn toegestaan:
\begin{prooftree}
  \AxiomC{$\lexists[x][!A(x)]$}
  \AxiomC{$\Discharge{!A(a)}{1}$}
  \RightLabel{*\Intro{\lforall}}
  \UnaryInfC{$\lforall[x][!A(x)]$}
  \RightLabel{\Elim{\lexists}}
  \BinaryInfC{$\lforall[x][!A(x)]$}
\end{prooftree}
Maar $\lexists[x][!A(x)] \Entails/ \lforall[x][!A(x)]$.

Omdat de eliminatieregels voor kwantoren alleen toestaan dat gesloten
termen voor !!{variable}s worden gesubstitueerd, volgt dat elke
!!{formula} die uit een verzameling !!{sentence}s kan worden afgeleid,
zelf !!a{sentence} is.
\end{explain}

\end{document}
